% Generated draft: requires mathematical and source review.
\documentclass[11pt]{article}
\usepackage[margin=25mm]{geometry}
\usepackage{fontspec}
\setmainfont{Latin Modern Roman}
\usepackage{amsmath,amssymb,mathtools,unicode-math}
\setmathfont{Latin Modern Math}
\usepackage{enumitem,needspace,xurl,hyperref}
\setlength{\parindent}{0pt}
\setlength{\parskip}{5pt}
\setlength{\emergencystretch}{3em}
\newcommand{\N}{\mathbb{N}}
\newcommand{\Z}{\mathbb{Z}}
\newcommand{\Q}{\mathbb{Q}}
\newcommand{\R}{\mathbb{R}}
\newcommand{\C}{\mathbb{C}}
\newcommand{\F}{\mathbb{F}}
\newcommand{\A}{\mathbb{A}}
\newcommand{\PP}{\mathbb{P}}
\newcommand{\E}{\mathbb{E}}
\newcommand{\eps}{\varepsilon}
\newcommand{\dd}{\,\mathrm{d}}
\begin{document}
\section*{UnsolvedMath ID: 20002682}
\textbf{Source problem code:} catalogue entry
\subsection{Definitions and mathematical statement}
\textbf{Question 22, including its first sentence.} Find a combinatorial interpretation of free cumulants through enumeration and operations on maps. The same question block then asks about rank-one spherical integrals. Let \(E_n\) be self-adjoint \(n\times n\) matrices with uniformly bounded spectra and empirical laws tending weakly to \(\mu_E\). With unnormalized matrix trace and Haar probability measure on \(\mathcal U(n)\), put
\[
 D_n=\operatorname{diag}(z,0,\ldots,0),\qquad
 I_n(z)=\int e^{n\operatorname{Tr}(UD_nU^*E_n)}\,\mathrm dm_n(U),
 \qquad f_n(z)=\frac{I_n'(z)}{nI_n(z)}.
\]
The logarithmic derivative is meromorphic and equals \(\partial_z(n^{-1}\log I_n)\) on zero-free neighborhoods, without requiring a global logarithm. The source attributes convergence near zero to \(R_{\mu_E}\) to Guionnet and Ma\"ida. Here the \(R\)-transform is the germ \(G_{\mu_E}^{-1}(z)-1/z\), using the Cauchy-transform inverse near infinity/zero. Determine the largest complex convergence domain, ask whether the whole upper half-plane works when \(\mu_E\) is freely infinitely divisible, and ask for noncompact-support extensions. The source suggests normal-family methods. The trace normalization, convergence topology, zero avoidance and analytic branch must be fixed; weak empirical convergence alone does not formulate all these choices.

\textbf{Reported balanced two-atom family.} Restrict to even \(n\), with half the eigenvalues equal to \(a\), half to \(b\), \(a>b\). Set \(c=(a+b)/2\), \(d=(a-b)/2\), \(\nu=(n-1)/2\). The aid reports
\[
 I_n(z)=e^{ncz}\Gamma(\nu+1)
 \left(\frac2{ndz}\right)^\nu I_\nu(ndz),\qquad
 f_n(z)=c+d\frac{I_{\nu+1}(ndz)}{I_\nu(ndz)}.
\]
Here \(I_\nu\) is the modified Bessel function, and the normalized first expression is interpreted by its entire power series at zero; \(f_n(0)=c\) is removable. If \(j_{\nu,1}\) is the first positive zero of the ordinary Bessel function \(J_\nu\), the nearest poles occur at
\[
 z=\pm\frac{\mathrm i j_{\nu,1}}{nd}
 \longrightarrow\pm\frac{\mathrm i}{a-b}.
\]
The limiting transform is
\[
 R(z)=c+\frac{\sqrt{1+4d^2z^2}-1}{2z},
\]
with square root equal to one at zero and removable value there. The aid reports locally uniform holomorphic convergence on the natural cut plane
\[
 \C\setminus\{\mathrm i y:|y|\geq1/(a-b)\},
\]
and exact maximal centered convergence-disk radius \(1/(a-b)\). This is a sharp family-specific obstruction; no assertion that the two-atom law is freely infinitely divisible, or that it settles the infinitely-divisible upper-half-plane question, is made. Nor is a particular cut plane identified as the uniquely largest possible analytic-continuation domain.

\textbf{Reported sufficient normal-family criterion.} Write \(I_n(z)=\mathbb E e^{nzY_n}\) for the rank-one spectral average, with \(|Y_n|\leq K\) under the uniform spectral bound. A precise sufficient phase-noncancellation condition on a connected domain \(D\) is: for each compact \(C\subset D\), there is \(\eta_C>0\), independent of \(n\), such that
\[
 |I_n(z)|\geq\eta_C\mathbb E e^{n\operatorname{Re}(z)Y_n}
 \quad(z\in C).
\]
Then \(I_n\) is zero free there and \(|f_n(z)|\leq K/\eta_C\). If \(D\) contains the established small-zero germ, local boundedness gives normal-family continuation to the corresponding analytic limit. This is a conditional test; no verification on a general domain or substitute tail condition for noncompact spectra is supplied. It does not address the combinatorial map question.

\textbf{Appended separate section fragment.} The extraction continues into section 0.9, ``Focus Group on Invariant Subspaces (day 4)''. For commuting \(T_1,\ldots,T_n\) in a \(\mathrm{II}_1\) factor, it proposes, with an explicit question mark, a joint Brown-measure identity involving
\[
 \log\Delta\left(1-\sum_i\alpha_iT_i\right)
 \quad\text{and}\quad
 \int\log\left(1-\sum_i\alpha_i\zeta_i\right)\,\mathrm d\mu(\zeta).
\]
Here \(\Delta(A)=\exp\tau(\log|A|)\) is the Fuglede--Kadison determinant, possibly zero. The printed complex logarithm without absolute values is not directly a real logarithmic-potential identity; a candidate repair would use \(\log|1-\sum_i\alpha_i\zeta_i|\), subject to integrability and extended-value conventions. The fragment supplies no complete existence/uniqueness hypotheses or requested theorem. Commutativity alone is not silently strengthened to joint normality, and this incomplete neighboring discussion is not treated as part of the spherical-integral result.
\subsection{Short English statement}
The source combines a cumulant/map question with complex spherical-integral questions and an appended Brown-measure fragment. The aid solves one balanced two-atom family, where Bessel poles approach exact \(R\)-transform branch points, and gives a conditional noncancellation test for normal-family continuation. The other questions remain separately qualified.
\subsection{Sources}
\begin{itemize}
\item Local UnsolvedMath registry, permanent ID \(20002682\), catalogue entry, \texttt{aim-probability-notes.json}, \texttt{notes[123]}, \emph{Free analysis}, Question 22 and appended section 0.9: \url{https://aimath.org/WWN/freeanalysis/freeanalysis.pdf}.
\item Original attribution: Guionnet and Ma\"ida, small-complex-parameter rank-one spherical-integral convergence to the \(R\)-transform. No full bibliographic reference is supplied.
\item Supplied research summary on the balanced two-atom Bessel formula, nearest-pole barrier and phase-noncancellation criterion. Attempt 1: partial\_result, partially\_solved, rigor ``proved'', low novelty confidence. The machine-generated aid is not peer reviewed and requires independent verification before citation. The source is labelled reconstructed\_unverified; its verification reports PDF inspection, which this draft has not independently performed.
\end{itemize}
\end{document}
